\chapter{Localisation}
\label{chapter:localisation}

So far, we have assessed our code as one big black box.

\begin{itemize}
  \item We did not ask \emph{which part of the code} causes a problem. We call
  these the \emph{hotspot}s of a code.
  \item We did not ask \emph{when} it arises throughout the execution.
  \item We did not ask \emph{where}, i.e.~on which rank (if we run over multiple
  nodes) a problem arises.
\end{itemize}

\noindent
Before we continue to assess the code and dive deeper into details, it makes
sense to gather the information from the three bullet point above or at least to
discuss variants how to gather it.
Hereby, we have to take the type of flaw that we are after into account:
Single core flaws imply that the ``where'' question does not play a role.
Inter-node flaws might suggest that the ``where'' question is the number one
question to ask.
As so often, all depends on its context.
We need the right tool for the right purpose.


\begin{assessmentcrime}
 80\% of the runtime might be spent on 20\% of the code.
 For the performance analysis, it is not appropriate to focus on this 20\% of
 the code only.
 Instead, you have to focus on all code that does not do any science if this
 code share is too big.
\end{assessmentcrime}

Too big means more than 20\%.
